\section{Guenon: The Trajectory of his Life}

For anyone looking for clues of an initiatic tradition in the West, it is helpful to follow the path trod by \textbf{Rene Guenon}. We can leave out of consideration any rumours or allegations, such as a mysterious Hindu or Sufi initiation. His own course speaks for itself and it is absurd to consider it as some sort of charade. The paths he actually followed are what counts, not the paths he may have studied. Of the former, there are three.

\paragraph{Hermetism}

At the end of the 19\textsuperscript{th} century, Paris was undergoing a sort of Hermetic revival, which attracted the young Guenon. He was secretary to Papus, one of the primary figures of that revival. He later abandoned that movement of some doctrinal issues and the lack of a regular initiation. However, keep in mind that this is the first place Guenon went to look for it.

\paragraph{Catholicism}


He next got involved with Catholic circles, even marrying a devout woman. He wrote for the journal Regnabit (``he Christ will reign'') and traveled in a circle that included the Thomist metaphysician \textbf{Jacques Maritain}. He supported the organization Action Française, led by Charles Maurras, which included Positivists and Catholics among its members. Interestingly, he supported the condemnation of Action Française by \textbf{Pope Pius XI} on the grounds he was the legitimate spiritual authority.

\paragraph{Sufism}


Unsatisfied with his first two paths, Guenon took advantage of a trip to Egypt following the death of his first wife to become fully engaged with Sufism. As we have pointed out, this was not a conversion in the conventional sense of the term\footnote{\url{http://www.gornahoor.net/?p=262}}. It involved an intellectual conversion and the personal judgment that Sufism provided a more certain initiatic for his needs. From an esoteric perspective, the abyss from Europe to Egypt may not be as wide as some think. \textbf{Houston Smith} describes Islam as a Semitic Religion built on Platonic metaphysics; he may as well have been describing Medieval Christendom. Furthermore, Sufism also has a Hermetic element of its own.

\paragraph{The Next Step}


Let's follow this logic, at least those of us who claim to be following the path created by Guenon. If the obvious places to look --- Hermetism and Medieval Catholicism --- are lacking a contemporary initiatic organization, then how can we expect to recognize such a thing? These obvious places are close to us in time, language, and history and are well documented. Presumably, they should be easier to understand than more alien texts such as the Yoga Sutras of \textbf{Patanjali}. Guenon mentions that the Romance of the Rose is an initiatic text, something he should be able to recognize.

Shouldn't all so-called Traditionalists rush out to obtain and interpret this work? I doubt it, as many will prefer to learn the alleged esoteric secrets of the Runes. I say alleged not because the Runes had no such significance, but only because no one today has any means to know those secrets.

What are the clues we can look for? First of all, when Reason speaks, this is understood as the higher intellect, or gnosis, something few contemporary minds can even conceive of. The book is ostensibly about profane love, but that clearly represents a description of a spiritual path, from \textbf{Dante} and the \textbf{Fedeli d'Amore} to the Sufi poets. A forgettable neo-pagan web site I came across suggested that the Medievals ``invented'' Mary because they had no proper goddesses of their own. That is merely a confession that they invented their neo-pagan goddeses.

Of course, the Medievals knew about \textbf{Sophia} and the \emph{Anima Mundi}, unlike those neo-pagans. The Romance even discusses the gods, a page they probably missed. They then make the tired claim that only some neo-pagan witches or Druids can understand it, while simultaneously misunderstanding it. To the contrary, Guenon writes about:

\begin{quotex}
Christian Hermetists of the Middle Ages, whose intentions were always fully orthodox.
\end{quotex}


Further East, and this is getting to the point, Guenon recognizes that

\begin{quotex}
Hesychasm whose truly initiatic character seems indisputable. ... In hesychasm, initiation in the strict sense consists essentially in the regular transmission of certain formulas, exactly comparable to the transmission of mantras in the Hindu tradition and of the wird in the Islamic turuq. It also contains a complete technique of invocations as a true method of interior work.
\end{quotex}


Guenon mentions that this tradition probably was able to continue by hiding itself. Shouldn't we recognize, however, in the writings of \textbf{Vladimir Solovyov}, \textbf{Valentin Tomberg}, and \textbf{Boris Mouravieff}, all coming from the East, an unveiling of that tradition. The elements are there: invocations, interior work, gnosis, Sophia and so on. I know there are those driven by sentimentalism and an overactive subjective imagination who would choke on such a thought. However, like the Aghoris who eat the repulsive, I say Chow down. A man on the path will stop at nothing in search of enlightenment.


\flrightit{Posted on 2011-12-12 by Cologero}

\begin{center}* * *\end{center}

\begin{footnotesize}
\begin{sffamily}
\texttt{Boreas on 2012-01-13 at 05:56 said:}

``Replacing one exoterism with an allegedly older, non-existent, poorly understood, and artificially constructed exoterism (neo-paganism) is a task for fools.''

It is not only a task for fools, it is mostly either disguised (or open) vanity, and, more importantly, it is to feed upon dead corpses. Maybe a single ``aghori'' here and there can do this and even manage to get some heavenly bread for one's self from it, but many either fool themselves and others at the same time and / or suffer the results of eating in the graveyard: a spiritual diarrhea and possibly even the curses of the dead upon one's souls.

``Let the dead bury the dead''... ?

\hfill

\texttt{Boreas on 2012-01-13 at 06:03 said:}

... and as a continuation to the above-mentioned (and my former posts

in this topic), the revival and serious study of ancient esoteric knowledge and esoteric traditions of one's own people, for example -- which is very understandable and possibly at the very basis of the yearnings of the neo-pagans also -- is a completely different thing.

This is to return to Tradition with a capital T, which is the same thing as to achieve a true palingenesis and a re-connection to the World Tree / Tree of Eternal Life through which veins the golden streams of love and light surge with ever-renewing Vigor, and not just one of its fallen and / or dried branches.

\hfill

\texttt{Charlotte Cowell on 2012-01-15 at 08:28 said:}

do we really have to lump feminisation together with vegetative?! esotericism is intimately concerned with the divine feminine so it seems counter productive (even counter-initiatory!) to negate 50 percent of God and humankind.

\hfill

\texttt{Boreas on 2012-01-15 at 14:11 said:}

observer: It seems that I have completely missed your somewhat lengthy, thoughtful and erudite comments about the guidestone issue until I now started to browse through the comments of this particular article.

All I can really say that I can relate to most your thoughts about the issue. Please notice that I took the ten points as themselves, and intentionally passed by the consideration of other issues related to it, which can be true or not. The ten points are somewhat ambivalent towards their interpretation, that's my point; they can be seen as subversive and they can be seen otherwise. They (and Freemasonry, NWO etc.) don't belong to my interests, nor do I find it useful to delve upon the matter anymore, after studying it from every possible angle for years and almost scorching myself totally because of it; this is precisely what the counter-initiation wants, for one to get involved and TOO interested about the world process and the course of history (`men in time' etc.)! Without a second doubt there are many shady things and plots going on in the world, and you can also take my word that I'm no beginner in this issue, that's for certain.

As for ``fighting the NWO'' part and becoming it's worker in the process, there's a good article about it in Cakravartin by Victor:

\url{http://www.cakravartin.com/archives/the-mother-of-all-conspiracies}


Evola himself saw it best to withdraw into seclusion and ``not to fight against evil'' (Ride the Tiger, the first essay) since the world has gone too far in the Kali-Yuga and the ``powers that be'' right now are on a short leash despite their apparent force, power and the ability for control. This was my point in that it is best to refrain from ``fighting the NWO'' (besides the great points made by Victor in the Cakravartin text), since it would be to invite a disaster both to one's self and to the things we (the men and women of Primordial Tradition) hold dear; look up for the article `A Way Out, Further In' here in Gornahoor. Tradition cannot be erased, since it is beyond time and the course of history, and it will prevail in the end no matter what.

Salve Victoria! Salve Nike! (Just do it?)

\hfill

\texttt{Boreas on 2012-01-15 at 14:50 said:}

A few more thoughts spanned to my mind about the issue. Logres seems to mention Rosenstock-Huessy here and there in his articles, and I became interested that has anyone read his book `Planetary Service -- A Way into the Third Millenium' or something like that? The man himself is a complete stranger to me and after a very short glance into the book at Google books it seems to point to somewhat same direction than the guidestones principles; can't be sure about this, however, since I have not read the book and know nothing about the man. However, he does seem to be valued by esotericists and ``traditionalists'' also.

As for the issue of ``a world state'', I think that looking the issue for once ``historically'', it is unavoidable in some form or another: how could anyone think that in this age of information technology etc., which is itself an ``exoteric'' sign of occult influence towards `a unification of humanity' in this or that form, could be avoided? I'm not trying to subvert anything here or impose my own agenda / ideology here, I am simply stating a fact. (Of course the real unification and a healthy `world state' would be altogether different than the one you have brought out, observer, with your comments about McDonalds, Monsanto, cultural levelling / marxism, ``one-world capitalist globalism'' etc. I certainly do not stand for any of these.)

As I would like to see it, it would be a world with healthy national cultures and true diversity, ruled by different-sized power blocks working in unison for the same cause, the cause of the `immortal dreamer' itself, not the current multicultural monolithic-globalist mish-mash of ``neutered zombies''. Nature we must certainly rever and honour, but you are absolutely right that it should not be placed above the spiritual cause.

Dixi.

\hfill

\texttt{Charlotte Cowell on 2012-01-15 at 16:36 said:}

well the antitode to the unhealthy zombified state of `monotone' socieites and so on is retreat to the heart of nature, that's the whole point. Spiritual focus is all very well but if we destroy what we was left in our charge in the process does that not in itself point to the adversary? pure spiritual impulses foster compassion towards all life, the world of creation. We must have our houses in order, both the temples of our souls but also the `Elysium fields' as they appear here. Do you not think that widespread pollution of water (lakes, rivers, etc), is symptomatic of general soul erosion, forgetfulness, greed, selfishness, etc? lack of mastery at any rate. As it is below so shall it be above. Same goes for animals -- we can learn so much from them at a soul/spirit level but they're held in contempt and abused by so many -- if we can't take care of innocent creatures how can we be entrusted with human souls and the development of humankind?

\hfill

\texttt{Charlotte Cowell on 2012-01-15 at 16:48 said:}

seeing as we have got onto the subject of the `occult war' (and I DO think it is within topic, though you are free to differ, of course, given that it is to do with counter-initiation as part of this), I am grateful to at least have it acknowledged. I have struggled to get true esoteric adepts to engage with me on how best to tackle the situation, which to me was an imperative duty of the call, part of graduating from school, so to speak. It seems to me that there is an intensified oppression present in certain spheres that is remarkably intense and demands constant vigilance. Of course vigilance is key at any stage, but my concern is also related to recent comments I made to do with harmonics and spiritual protection, such as fortifying the aura, donning the armour of God in full, etc. Particularly if one's guardian angel has been `released' into full service to the spiritual hierarchy above. Never have we been more in need of Michael. I'm sure most of you would do this naturally anyway, but I think it is worth pointing this out in case anyone is dozing... .thief in the night

\hfill

\texttt{Charlotte Cowell on 2012-01-15 at 16:52 said:}

I mean surely, for example, it's time for spiritual hostages to be freed (while of course we individually clear through what purgatory we can)? I think it is time for this to happen, for sleepers to arise, including myself in many aspects! do others not?

\hfill

\texttt{logres on 2012-01-15 at 22:20 said:}

Women \& men approach the same objects very differently. Charlotte, you might check out William Humboldt's letters (in Humanist Without Portfolio) on the spirituality of women -- I found them helpful, and they seemed to accord with what I've experienced, in interacting with the fair sex. If I remember correctly, he claims women have a harder time actualizing, but if they do, they make a more comprehensive approach to spiritual matters. I've been caught up in the NWO order thing too much myself -- Boreas' link was very helpful (thanks). But your point about Nature \& animals is timely and not really debateable (in my mind), and I certainly don't disagree that some awareness of all these currents is helpful in our growth.

\hfill

\texttt{Boreas on 2012-01-16 at 04:33 said:}

Charlotte is very correct that our relationship towards earth, nature and animals reflects the state of our souls and our relationship towards the divine feminine itself. Evola and many traditional authors have had quite a dangerous relationship towards `The Mother'. For example, Evola refused the attempts of healing and miracles from a catholic priest who wished to heal him in the name of the Holy Mother -- yet he lived with his mother for a long time. Quite telling, isn't it!? Too much solarity leads easily to scorching (Evola was wounded in the Manipura chakra), which makes Evola very `combustible', as has been pointed out here. Masculinity and solarity should always be balanced with feminity and the right kind of lunarity -- the alchemical conjuction of the Sun and Moon etc.

I think these books and authors can do alot to clear the issue between spirituality, tradition and ecology.

\url{http://www.sophiaperennis.com/category/books/spiritual-ecology/}


I have myself read only the `Looking Back on Progress' by Lord Northbourne, and it was a true eye-opener also towards many other things besides ecological questions, but what I have thus far read from Charles Upton, has also convinced me of his humble sufi wisdom.

The exoteric monotheisms have been very destructive in this regard, as Mr. Upton shows already in the presentation of his book `Who Is the Earth'. It could also be said that the process from the sphere to the cube in its negative repercussions underlies the weltanshauung of all the totalitarian systems.

\hfill

\texttt{Charlotte Cowell on 2012-01-16 at 11:22 said:}

Oddly enough (while I realise it's considered a masculine trait) up until a couple of years ago I had an enormous excess of fire, which led to all sorts of difficulties and ultimately a form of burnout, although it also helped to give me the spiritual vision and second sight. Creation of fiery water from this state led to the floodgates being opened

\hfill

\texttt{Boreas on 2012-01-16 at 12:11 said:}

For the sake of clarity I need to clear one thing from the comment 55, since it could be easily interpreted that I'm lying. That is, how could I know the R-H's book's name if I know absolutely nothing about the man? I glanced quickly through his wikipedia biography once earlier on, when I saw Logres mentioning his name here, and the book name caught my eye then already, and now it came back to my mind when this issue popped out.

To observer I wish to say that hopefully you have not been irritated by my sometimes irritable rhetorics. I live in a mad mansion with small children and occult paint and oil vapors, and I cannot always be held 100 \% responsible of what kind of frogs come out from my mouth. I'm glad you questioned me and it's a good thing the issue was dealt with. Blessings of the High One to you.

\hfill

\texttt{Boreas on 2012-01-21 at 12:00 said:}

Although I said earlier on that the issue and argument is already finished for my part, I wish to get into the `guidestones' / masonry issue shortly, and this is also a reference -- and a deep veneration, for sure -- to Guénon and to his life's work. Although I'm generally well aware about the many times subversive and negative role that modern Freemasonry has played in the course of world history, Prompted by this I actually ordered Guénon's book about Freemasonry and related issues to see his thoughts upon the matter. I have not yet read the book but only got a glance to it here and there, and that was certainly enough to realize also the guidestones issue quite throughly: instead of `the Great Work of Universal Construction' which is the purpose of \_ancient and true masonry\_, it is the manifestation of the Tower of Babel and ``the confusion of tongues'' only makes it more blatanly clear. (I don't know if observer is still reading the topic and this forum, but kudos for you.)

\hfill

\end{sffamily}
\end{footnotesize}

